% 第 5 讲，源文件 src/lecture05.txt（317 行）；R3 块。
\section*{第 5 讲\quad 阿代尔配对的非退化性与正合列}

回忆一下，我们想要一个非退化的配对，而
\begin{equation}\tag{5.1}
\Q^\times\backslash\mathbb{A}_\Q^\times/(\mathbb{A}_\Q^\times)^2\times\Q^\times/(\Q^\times)^2\longrightarrow\{1,-1\},\qquad
((x_p),r)\longmapsto\prod_p (x_p,r)_p
\end{equation}
曾是候选者（与之前一样，$p$ 取遍所有素数以及 $\infty$）。这个配对的良定义性归结为互反律
\[
\prod_p (x,y)_p=1\qquad (x,y\in\Q^\times).
\]
我们看到，当 $x=p$、$y=q$ 为奇素数时，它归结为二次互反律
\[
\left(\frac{p}{q}\right)\left(\frac{q}{p}\right)=(-1)^{\frac{p-1}{2}\cdot\frac{q-1}{2}},
\]
而我们是通过考虑特征标
\[
\chi\colon\Gal(\Q(\zeta_p)/\Q)=(\Z/p\Z)^\times\xrightarrow{\ (\cdot/p)\ }\{1,-1\}
\]
来证明它的。我们看到，这对应于 $\Q(\zeta_p)$ 的一个唯一二次子扩张
\[
\Q(\sqrt{\pm p})=\Q\Bigl(\sqrt{\left(\tfrac{-1}{p}\right)p}\Bigr),
\]
其中的要点是高斯和
\[
G:=\sum_{a=1}^{p-1}\left(\frac{a}{p}\right)\zeta_p^a=\sqrt{\left(\frac{-1}{p}\right)p}.
\]

更一般地，若 $F/\Q$ 是 Galois 扩张、Galois 群为 $G$，且 $\Q$ 的一个素位 $q$ 非分歧，则 $[\mathrm{Frob}_q]\mapsto[1]\in G$（这里两边都是共轭类）当且仅当 $q$ 在 $F$ 中分裂。于是
\[
\left(\frac{q}{p}\right)=1\iff
\left(\frac{\left(\tfrac{-1}{p}\right)p}{q}\right)=1,
\]
因为右端等价于 $q$ 在扩张 $\Q(\sqrt{\pm p})$ 中分裂，而这又蕴含
\[
\left(\frac{q}{p}\right)=\left(\frac{\left(\tfrac{-1}{p}\right)p}{q}\right)
=\left(\frac{\left(\tfrac{-1}{p}\right)}{q}\right)\left(\frac{p}{q}\right)
=(-1)^{\frac{p-1}{2}\cdot\frac{q-1}{2}}\left(\frac{p}{q}\right),
\]
由此即得所需结论。

类似地，我们可以在其他情形下得到互反律的相应结论，例如：

\textbf{命题 5.1.} 我们有
\[
\prod_p (2,\ell)_p=1,
\]
其中 $\ell$ 为奇素数，$p$ 取遍所有素数（注意 $(2,\ell)_\infty=1$ 是平凡的）。

\textit{证明.} 与之前一样，若 $p\neq 2,\ell$ 则 $(2,\ell)_p=1$；而利用驯符号，
\[
(2,\ell)_\ell=\left(\frac{(-1)^{v(\ell)v(2)}2^{v(\ell)}/\ell^{v(2)}}{\ell}\right)
=\left(\frac{2}{\ell}\right).
\]
由习题集 1 第 2(d) 题所得的公式，我们有
% [已核：原文该式指数中原有两处在 PDF 抽取里成了不可识别字符（U+000F），按原 PDF 第 1177 行（Serre《A Course in Arithmetic》第 III 章的 Q_2 公式）读作 ε(1)ε(ℓ)+v(2)θ(ℓ)+v(ℓ)θ(1)；
%  由于 ε(1)=0、v(ℓ)=0，该式化为 (-1)^{θ(ℓ)}，与右端结论一致]
\[
(2,\ell)_2=(-1)^{\varepsilon(1)\varepsilon(\ell)+v(2)\theta(\ell)+v(\ell)\theta(1)}=(-1)^{\theta(\ell)},
\]
其中
\[
(-1)^{\theta(\ell)}:=
\begin{cases}
1, & \ell\equiv 1,-1\pmod 8,\\
-1, & \ell\equiv 3,-3\pmod 8,
\end{cases}
\]
这对应于 $\Z/16\Z$ 中 $\ell^2$ 到 $\Z/2\Z$ 的标准同构。因此，我们希望证明
\[
\left(\frac{2}{\ell}\right)=(-1)^{\theta(\ell)}.
\]
为了判断 $2$ 是否为模 $\ell$ 的平方，我们想要 $\sqrt2$ 的一个便利表达式，也就是 $\Q(\sqrt2)$ 的一个分圆嵌入（在其中 $\ell$ 分裂当且仅当 $\left(\frac{2}{\ell}\right)=1$）。回忆若 $\zeta_8$ 为本原八次单位根，则可取 $\zeta_8=\frac{\sqrt2}{2}+i\frac{\sqrt2}{2}$，于是
\[
\zeta_8+\zeta_8^{-1}=\zeta_8+\overline{\zeta}_8=2\operatorname{Re}(\zeta_8)=\sqrt2.
\]
在代数上，可以这样看出这一恒等式：
\[
(\zeta_8+\zeta_8^{-1})^2=2+\zeta_8^2+\zeta_8^{-2}=2+\zeta_4+\zeta_4^{-1}=2,
\]
因为 $\zeta_4$ 与 $\zeta_4^{-1}$ 恰为 $i$ 与 $-i$。这给出 $\Q(\sqrt2)\subseteq\Q(\zeta_8)$，以及特征标
\[
\Gal(\Q(\zeta_8)/\Q)=(\Z/8\Z)^\times\xrightarrow{\ \chi\ }\{1,-1\}.
\]
我们断言 $\left(\frac{2}{\cdot}\right)=\chi=(-1)^{\theta(\cdot)}$。显然 $\ker\bigl((-1)^{\theta(\cdot)}\bigr)=\{1,-1\}$，而 $n\in(\Z/8\Z)^\times$ 属于 $\ker(\chi)$ 当且仅当它固定 $\sqrt2=\zeta_8+\zeta_8^{-1}$，即
\[
\zeta_8+\zeta_8^{-1}\longmapsto\zeta_8^n+\zeta_8^{-n}=\zeta_8+\zeta_8^{-1},
\]
而这仅在 $n=1$（两项都不动）或 $n=-1$（两项互换）时成立。于是两个核相同，从而两个函数相等。\hfill$\square$

注意，我们也可以不用 Galois 群来论证。若假定 $\zeta_8+\zeta_8^{-1}\in\F_\ell$，那么要验证这一点，只需检查
\[
\zeta_8+\zeta_8^{-1}=(\zeta_8+\zeta_8^{-1})^\ell=\zeta_8^\ell+\zeta_8^{-\ell},
\]
也就是说它在 Frobenius 元的作用下不动；于是得到与前面相同的条件。

其他符号验证起来相对繁琐：例如 $\prod_p(2,2)_p=1$ 很简单，因为由习题集 1 第 2(d) 题有 $(2,2)_2=1$；而 $\prod_p(-1,\ell)_p=1$ 则由
\[
\left(\frac{-1}{p}\right)=(-1)^{\frac{p-1}{2}}
\]
看出。至此我们已验证了初始配对 (5.1) 的良定义性。现在我们要验证这个配对是非退化的。注意这一步其实并不需要互反律，因为论证反而更容易。

\textbf{命题 5.2.} (5.1) 中定义的映射
\[
\chi\colon\Q^\times\backslash\mathbb{A}_\Q^\times/(\mathbb{A}_\Q^\times)^2
\xrightarrow{\ \sim\ }\operatorname{Hom}(\Q^\times,\{1,-1\})\cong\Gal_2(\Q)
\]
是同构（注意：配对只从 $\Q^\times/(\Q^\times)^2$ 出发定义映射并无妨碍，因为右端的同态会吸收平方）。

\textit{证明.} 在习题集 2 第 1(a) 题中，我们证明了
\[
\mathbb{A}_\Q^\times=\widehat{\Z}^\times\times\R_{>0}\times\Q^\times,
\]
其中前两项经由局部素位嵌入，最后一项沿对角线嵌入。模去 $\Q^\times$ 消去最后一项，模去平方消去第二项，于是由中国剩余定理得到
\begin{equation}\tag{5.2}
\Q^\times\backslash\mathbb{A}_\Q^\times/(\mathbb{A}_\Q^\times)^2=\prod_p\Z_p^\times/(\Z_p^\times)^2,
\end{equation}
其中 $p$ 取遍所有素数；当 $p$ 为奇素数时，$\Z_p^\times/(\Z_p^\times)^2$ 是 2 阶群，由任一二次非剩余生成；而 $\Z_2^\times/(\Z_2^\times)^2\cong(\Z/8\Z)^\times$ 是 4 阶群，由 $-1$ 与 $5$ 生成。又有
\[
\Q^\times/(\Q^\times)^2=\{\pm p_1\cdots p_r:\ p_i\ \text{为素数}\}
=\Z/2\Z\oplus\bigoplus_p\Z/2\Z,
\]
其中 $p$ 同上，第一个 $\Z/2\Z$ 对应符号。我们将看到，在对偶之后，这些 $\Z/2\Z$ 的副本全都（几乎）一一对上。

设 $p$ 为奇素数，$r$ 为 $p$ 处的二次非剩余，即 $\Z_p^\times/(\Z_p^\times)^2$ 中的非平凡元。则
\[
(\chi(r)(\pm q))_p=(r,\pm q)_p
=\left(\frac{r^{v(\pm q)}/(\pm q)^{v(r)}}{p}\right)
=\begin{cases}1, & q\neq p,\\ -1, & q=p,\end{cases}
\]
是 $\chi$ 在 $\Q^\times\backslash\mathbb{A}_\Q^\times/(\mathbb{A}_\Q^\times)^2$ 的第 $p$ 个分量上所取的值，其中 $q$ 为素数。对最后一个基元素，有 $\chi(r)(-1)=1=(r,-1)_p$。于是 $p$ 处那些显然的（拓扑）基元素一一对上；现在我们必须问 $p=2$ 时会发生什么。一个自然的猜测是「在 $2$ 处取 $r=-1$ 或 $r=5$、其余各处取 $r=1$」所定义的伊代尔，因为由习题集 2 第 2(b) 题，$5$ 对应 $\Q_2$ 的唯一非分歧二次扩张。计算可得
\[
\begin{aligned}
\chi(5,1,1,\dots)(q)&=(5,q)_2=1,\\
\chi(5,1,1,\dots)(-1)&=(5,-1)_2=1,\\
\chi(5,1,1,\dots)(2)&=(5,2)_2=(-1)^{\theta(5)}=-1,
\end{aligned}
\]
可见这个基元素与 $2$ 处的基元素完美对应。这里我们用元组表示伊代尔，其坐标是就同构 (5.2) 而言取的，素数按通常顺序排列。于是
\[
\chi(-1,1,1,\dots)(-1)=-1=(-1,-1)_2
\]
完成了证明。不过，这两组基实际上并不完美对应，因为与另一个奇素数 $p$ 配对会给出与 $-1$ 是否为模 $p$ 的平方相应的符号；但我们可以用「上三角矩阵」轻易地把一组基用另一组表示出来，办法是对 $(-1,1,1,\dots)$ 这个基元素作修正（本质上，我们有一个对角线上全为 $1$ 的无限矩阵，只是在 $(-1,1,1,\dots)$ 处例外，它对应于 $\Q^\times/(\Q^\times)^2$ 的给定基中一个更复杂的元素）。

下面给一个稍微更严肃的论证。我们有如下行正合的交换图：
\[
\begin{CD}
0 @>>> \Z/2\Z\times\prod_{p\neq2}\Z/2\Z @>>>
\Q^\times\backslash\mathbb{A}_\Q^\times/(\mathbb{A}_\Q^\times)^2 @>>> \Z/2\Z @>>> 0\\[2.2ex]
@. @V{\sim}VV @V{\sim}VV @V{\varphi\mapsto\varphi(-1)}VV @.\\[2.2ex]
0 @>>> \{\varphi\colon\Q^\times\to\Z/2\Z\mid\varphi(-1)=1\} @>>>
\operatorname{Hom}(\Q^\times,\Z/2\Z) @>>> \Z/2\Z @>>> 0
\end{CD}
\]
第一个 $\Z/2\Z$ 对应伊代尔 $(5,1,1,\dots)$，其余的副本对应每个 $p$ 处的二次非剩余；在最右端的 $\Z/2\Z$ 中，我们得到 $(-1,1,1,\dots)$ 的像。两个映入 $\Z/2\Z$ 的映射都是商映射；右侧的竖直映射是同构，因为它非平凡；左侧的竖直映射也是同构，因为如我们先前所见，一切完美对应。于是中间的映射是同构，此即所求。\hfill$\square$

现在回到这个问题：证明对局部域的任何二次扩张 $L/K$，都有 $\#(K^\times/NL^\times)=2$。回忆这一陈述等价于希尔伯特符号的双乘性与非退化性；我们已经对奇素数以及非分歧、驯分歧扩张的情形证明了它，但对野分歧扩张、以及对 $\Q_2$ 上除 $\Q_2$ 本身之外的扩张还证不出来。我们现在的目标是更一般地证明它：即证明若 $L/K$ 是 $n$ 次循环扩张（也就是说，它是 Galois 群为 $\Z/n\Z$ 的 Galois 扩张），则 $\#(K^\times/NL^\times)=n$。为了把它放进更一般的框架：若 $L/K$ 是有限阿贝尔扩张，我们实际上期望有典范同构
\[
\Gal(L/K)\cong K^\times/NL^\times,
\]
所以期望的比数值相等更多。我们将用正合列与同调代数的方法来证明这一点，现在转入这一话题。

事实证明，短正合列是确定有限阿贝尔群阶数的极好工具。设我们有短正合列
\[
0\to M\overset{g}{\hookrightarrow}E\overset{f}{\twoheadrightarrow}N\to 0.
\]
则 $M=\ker(f)$，$N=\operatorname{coker}(g)=E/M$；用滤链的语言说，$M$ 与 $N$ 就像相伴分次项。确实，可以把 $M$ 与 $N$ 看作「原子」，把 $E$ 看作「分子」，而分子的精细结构决定它的「反应」。显然若 $M$ 与 $N$ 有限，则 $E$ 也有限，且 $\#E=\#M\#N$。野分歧带来的问题在于：$L^\times$ 上没有滤链。

一个问题是，许多运算并不保持短正合列。例如，若 $n\ge1$ 为整数，模去 $n$ 并不保持 $\#(E/nE)$。

\textbf{例 5.3.} (1) 若我们有正合列
\[
0\to\Z/n\Z\xrightarrow{x\mapsto(x,0)}\Z/n\Z\times\Z/n\Z
\xrightarrow{(x,y)\mapsto y}\Z/n\Z\to 0,
\]
则模去 $n$ 保持它。

(2) 若我们有正合列
\[
0\to\Z/n\Z\xrightarrow{x\mapsto xn}\Z/n^2\Z\xrightarrow{1\mapsto1}\Z/n\Z\to 0,
\]
则模去 $n$ 把它变成
\[
\Z/n\Z\to\Z/n\Z\xrightarrow{\operatorname{id}}\Z/n\Z\to 0.
\]
我们有同样的「原子」，但它们构成了不同的「分子」。在前一情形，模去 $n$ 之后的阶是 $n^2$，而在这里是 $n$。

同调代数的一个中心论点是：我们可以通过延长正合列来修正这一点。打个诗意的比方，被改变的正合列就像可见光；它缺了红外波段，而延长正合列就能让我们看到这一段。具体说来，这对应于 $n$-挠：$\Z/n^2\Z$ 没有 $\Z/n\Z\times\Z/n\Z$ 那么多的 $n$-挠，后者全是 $n$-挠。

对模 $M$，记 $M[n]\subseteq M$ 为 $\ker(n\colon M\to M)$。于是我们得到更长的正合列
\[
0\to M[n]\to E[n]\to N[n]\xrightarrow{\ \delta\ }M/n\to E/n\to N/n\to 0,
\]
其中对 $x\in N[n]$，$\delta(x)=ny$，这里 $y\in E$ 是满足 $f(y)=x$ 的任意元素；注意 $f(ny)=nx=0$，故 $\delta(x)\in M\subseteq E$，此即所求。我们将在下一讲证明这是一个正合列。
